\section{Position effects and compositional difficulty}\label{sec:C11:geometry}
Length is only one coordinate of task difficulty. A useful workload specification includes evidence distance, evidence density, distractor similarity, the number of required joins, contradiction depth, and the delay between observation and consequence. Evidence density is the fraction of input that contributes to a correct answer. Low density stresses selection; high density with many dependencies stresses integration. Two inputs of equal token length can therefore create different computational demands.

Lost in the Middle provides evidence that information placement can affect model performance in its tested retrieval and question-answering settings~\cite{liu2023lostmiddle}. RULER and LongBench broaden long-context evaluation through controlled and heterogeneous tasks~\cite{hsieh2024ruler,bai2023longbench}. Their central methodological value here is to make context use observable under specified tasks. They do not establish a universal conversion from advertised token capacity to a reliable operating horizon. Synthetic retrieval, document question answering, and multi-step interactive execution expose overlapping but different failure modes.

The pinned RULER implementation makes several controls inspectable. In \path{scripts/data/synthetic/niah.py}, \path{generate_input_output} varies needle keys, values, queries, distractors, and insertion positions; \path{generate_samples} uses the selected tokenizer and reserves generated tokens and model-template overhead while fitting inputs. The standalone \path{scripts/pred/call_api.py} prediction path appends the stored answer prefix before dispatch. In \path{scripts/eval/evaluate.py}, \path{run_evaluation_per_task} applies the task metric and separately counts empty predictions~\cite{cumulative_a_ruler_repo_20261029}. Preserve the exact rendered input, tokenizer, seed, target position, and generation reserve when comparing lengths. The inspected prediction helper retries exceptions without a finite retry limit, and the evaluator skips absent task files after printing a warning. A reproducible adaptation should impose a declared deadline, reconcile expected and completed task identifiers, and retain failed requests in its accounting. These are implementation obligations; this static benchmark does not itself certify session continuity or correct external effects.

To separate these effects, construct a family of tasks with fixed decisive records and systematically permute their positions while keeping total length constant. Next, increase distractor similarity without changing the decisive records. Finally, change the number of relations needed to connect the records. The resulting response surface $s(L,d,k,\rho)$ records success as a function of length $L$, evidence distance $d$, join count $k$, and distractor parameter $\rho$. A length-only curve is a projection of this surface and can conceal strong interaction effects.

Repeated facts require special care. Repeating a constraint near the decision can improve accessibility while changing the workload's redundancy. If a long-context method benefits from explicit reminders, that result should be distinguished from recovery of a unique distant fact. Likewise, a task that includes both a stale fact and a corrected fact tests temporal interpretation as well as retrieval. An answerer can retrieve both perfectly and still resolve their authority incorrectly. Position permutations must preserve timestamps and source roles or the experimental manipulation changes the intended semantics.

Compositional difficulty also interacts with output format. Choosing a cited answer from a small candidate set can be easier than synthesizing an executable action with multiple arguments. An agent must often retrieve an entity, infer its current state, select a tool, satisfy a schema, and respect an authorization boundary. Evaluation that terminates after entity retrieval measures only the first component. For expert comparison, the action-generation stage should be retained, and failures should be attributed to evidence access, reasoning, argument construction, or execution.

An informative negative control inserts irrelevant but lexically attractive records whose instructions conflict with the user task. The model should recognize these as data rather than authority. This extends context evaluation toward the security analysis in Chapter~\ref{ch:C28}. Success under benign distractors does not imply robustness when retrieved content attempts to redefine the task or rewrite the memory policy.

